\usepackage[T1]{fontenc}
\usepackage{lmodern}
\usepackage{amsmath,amssymb,amsthm,mathtools,mathrsfs}
\usepackage[margin=1in]{geometry}
\usepackage{enumitem,microtype,booktabs,array,longtable}
\usepackage{xcolor,graphicx,tikz}
\usetikzlibrary{arrows.meta,positioning,calc,fit,decorations.pathreplacing,shapes.geometric}
\usepackage[hyphens]{url}
\usepackage[colorlinks=true,linkcolor=black,citecolor=blue!55!black,urlcolor=blue!55!black]{hyperref}
\usepackage[nameinlink,noabbrev,capitalize]{cleveref}
\hypersetup{pdftitle={The Poisson--Dirichlet law for prime predecessors},pdfauthor={OpenAI}}
\newcommand{\1}{\mathbf{1}}
\newcommand{\e}{\mathrm{e}}
\newcommand{\E}{\mathbb E}
\newcommand{\Prob}{\mathbb P}
\newcommand{\N}{\mathbb N}
\newcommand{\Z}{\mathbb Z}
\newcommand{\R}{\mathbb R}
\newcommand{\C}{\mathbb C}
\newcommand{\eps}{\varepsilon}
\newcommand{\dd}{\,\mathrm d}
\DeclarePairedDelimiter{\norm}{\lVert}{\rVert}
\DeclarePairedDelimiter{\abs}{\lvert}{\rvert}
\newcommand{\inner}[2]{\langle #1,#2\rangle}
\DeclareMathOperator{\Li}{Li}
\DeclareMathOperator{\supp}{supp}
\DeclareMathOperator{\cond}{cond}
\DeclareMathOperator{\lcm}{lcm}
\DeclareMathOperator{\rad}{rad}
\DeclareMathOperator{\PD}{PD}
\newtheorem{theorem}{Theorem}[section]
\newtheorem{lemma}[theorem]{Lemma}
\newtheorem{proposition}[theorem]{Proposition}
\newtheorem{corollary}[theorem]{Corollary}
\theoremstyle{definition}
\newtheorem{definition}[theorem]{Definition}
\theoremstyle{remark}
\newtheorem{remark}[theorem]{Remark}
\numberwithin{equation}{section}
\setlist[enumerate]{itemsep=3pt,topsep=5pt}
\setlist[itemize]{itemsep=3pt,topsep=5pt}
\allowdisplaybreaks[1]
\emergencystretch=1em
